\scanpage{286}
également un ``bon'' sous-schéma abélien.

\margin{$\unsure{\rho}$}
[N.B. \unsure{Il n'était} pas nécessaire dans \deleted{\illegible{1}} l'exemple plus général que $B$ soit un schéma abélien sur lequel $\Gamma$ opère \deleted{par}\inserted{par} translations.
En effet, en vertu de résultats généraux, on a une suite exacte
\diagram{\begin{tikzcd}[column sep=small]
0\arrow[r]&\deleted{\unsure{A'}}\inserted{\uuline{\operatorname{Pic}}_{A/S}}\arrow[r]&\uuline{\operatorname{Hom}}_S(B,\deleted{\unsure{A'}}\inserted{\uuline{\operatorname{Pic}}_{A/S}})\arrow[r]&\uuline{\operatorname{Hom}}_{S-\mathrm{gr}}(A,\uline{\operatorname{Pic}}_{B/S}^{\tau})\arrow[r]&0
\end{tikzcd}}
\margin{morphismes constants}
\margin{[correspond aux \unsure{bi-inversibles} sur $A\times B$]}
\margin{\unsure{sous les} conditions habituelles sur $B$\ldots}
d'où une suite exacte
\diagram{\begin{tikzcd}[column sep=small]
0\arrow[r]&\uuline{\operatorname{Pic}}_{A/S}^{\Gamma}\arrow[r]&\uuline{\operatorname{Hom}}_{S,\Gamma}(B,\uuline{\operatorname{Pic}}_{A/S})\arrow[r]&\uuline{\operatorname{Hom}}_{S-\mathrm{gr}}(A,\uuline{\operatorname{Pic}}_{B/S}^{\tau})^{\Gamma}
\end{tikzcd}}
Or \deleted{\unsure{tout}} supposons que $A_{\Gamma}=0$ \unsure{ensemblist!} \inserted{et $\Gamma$ \uline{opère trivial!} sur $\uline{\operatorname{Pic}}_{B/S}^{\tau}$}, alors
\[
\uuline{\operatorname{Hom}}_{S-\mathrm{gr}}(A,\uuline{\operatorname{Pic}}_{B/S}^{\tau})^{\Gamma}=0
\]
[utiliser le fait que c'est un schéma non ramifié sur $S$], et il vient
\[
\uuline{\operatorname{Pic}}_{A/S}^{\Gamma}\simeq\uuline{\operatorname{Hom}}_{S,\Gamma}(B,\uuline{\operatorname{Pic}}_{A/S}))
\]
Donc on aura
\[
\text{\textcircled{$*$}}\quad\left\{\begin{array}{l}
\uuline{\operatorname{Pic}}_{X/S}\simeq\uuline{\operatorname{Pic}}_{Y/S}\times\uuline{\operatorname{Pic}}_{A/S}^{\Gamma}\\[1.5em]
\text{d'où}\quad\uuline{\operatorname{Pic}}_{X/S}^{\tau}\simeq\uuline{\operatorname{Pic}}_{Y/S}^{\tau}\times A'^{\Gamma}
\end{array}\right.
\]
\deleted{]}
On peut remarquer qu'on peut supposer de \unsure{$=$}, au lieu que $A$ soit un schéma abélien, que $\operatorname{Alb}(A)$ existe, et $(\operatorname{Alb}(\unsure{M|A}))_{\Gamma}=0$ ensemblistiquement --.
]
\figure{\begin{tikzpicture}[x=1cm,y=1cm]
\draw (0,1.7)..controls (0.2,1.5) and (0.8,1.95)..(0.9,1.57)..controls (1.2,1.8) and (1.6,1.64)..(2,1.8);
\draw (0,0.7)..controls (-0.2,0.55) and (-0.2,0.4)..(0.1,0.35)--(2.5,0.3)--(2.55,0)..controls (2.6,0.38) and (2.65,0.3)..(3.2,0.3);
\end{tikzpicture}}
\mspage{285}
